% ECONOMICS_PROBLEM: {"id": "F37-573", "title": "Exchange-rate forecasting", "jel_code": "F37", "source_id": "OP-0330"}
% FIRST_POSTED: 2026-10-09
% ATTEMPT_STATUS: unresolved
% ENTRY_KIND: research_attempt
\documentclass[11pt]{article}
\usepackage{amsmath,amssymb,amsthm,hyperref}
\usepackage[margin=1in]{geometry}
\setlength{\emergencystretch}{2em}
\newtheorem{theorem}{Theorem}
\newtheorem{proposition}[theorem]{Proposition}
\newtheorem{lemma}[theorem]{Lemma}
\newtheorem{corollary}[theorem]{Corollary}
\theoremstyle{remark}
\newtheorem{remark}[theorem]{Remark}
\newtheorem{example}[theorem]{Example}
\title{Exchange-rate forecasting: predictable components and honest delayed comparisons}
\author{GPT 6 Astra Ultra}
\date{October 9, 2026}
\begin{document}
\maketitle

\section*{Target and scope}
For a prespecified currency pair and horizon $h$, the squared-loss target relative to the no-drift random walk is
\[
 \Delta_h(f)=E[(s_{t+h}-\widehat s^f_{t+h})^2-(s_{t+h}-s_t)^2].
\]
Negative loss difference is a statistical statement about a declared evaluation law. It is not equivalent to stationarity or profitable trading. Bakker's primary working paper \cite{exchange} explicitly combines permanent and transitory components; its conditional forecasting findings do not establish universal improvement. This note derives a transparent gain formula, quantifies predictor noise, and gives a sequential comparison allowing overlapping horizons.

\section*{An integrated process can be predictably forecastable}
Consider
\[
 s_t=R_t+z_t,\qquad R_{t+1}=R_t+\eta_{t+1},\qquad
 z_{t+1}=\phi z_t+\varepsilon_{t+1},\quad |\phi|<1.
\]
The two innovation sequences are mutually independent, centered, square integrable, and independent of the real-time past. Assume $z_t$ is stationary with variance $V=\sigma_\varepsilon^2/(1-\phi^2)$. All parameters below are supplied to isolate the forecasting question. Let the information set include $s_t$ and satisfy the stated future-innovation independence, and put $m_t=E[z_t\mid\mathcal I_t]$.

\begin{theorem}[The exact squared-loss gain]
The conditional-mean forecast is
\[
 s_t+(\phi^h-1)m_t.
\]
Its expected loss improvement over the random walk is
$(1-\phi^h)^2E[m_t^2]$. If $z_t$ itself is observed, its mean squared error is
\[
 h\sigma_\eta^2+\sigma_\varepsilon^2\sum_{j=0}^{h-1}\phi^{2j},
\]
whereas random-walk mean squared error adds $(1-\phi^h)^2V$.
\end{theorem}
\begin{proof}
Iteration gives
$s_{t+h}-s_t=(\phi^h-1)z_t+\sum_{j=1}^h\eta_{t+j}+\sum_{j=1}^h\phi^{h-j}\varepsilon_{t+j}$.
Conditional expectation annihilates future innovations and replaces $z_t$ by $m_t$. For any square-integrable outcome $U$, conditional-mean projection gives
$E(U-g)^2=E(U-E[U\mid\mathcal I_t])^2+E(g-E[U\mid\mathcal I_t])^2$
for every information-measurable forecast $g$. Apply this to the random-walk forecast to obtain the gain. When $z_t$ is observed, only the future innovation sums remain in the optimal error; independence gives the displayed variance. The random-walk error also contains $(\phi^h-1)z_t$, independent of those sums, adding the stated term.
\end{proof}
If the stationary component is absent or conditionally unpredictable, the gain is zero. In the pure random-walk submodel any forecast differing from $s_t$ has excess risk equal to the mean square of that difference. Thus a model class containing that submodel cannot promise strict improvement everywhere without further restrictions.

\begin{proposition}[Noisy predictors require shrinkage]
Suppose a supplied real-time predictor is $v_t=z_t+\xi_t$, where $\xi_t$ is centered, independent of $z_t$ and future innovations, with variance $\nu$. Consider forecasts
$\widehat s_{t+h}=s_t-\lambda(1-\phi^h)v_t$. Their loss difference is
\[
 \Delta_h(\lambda)=(1-\phi^h)^2\{\lambda^2(V+\nu)-2\lambda V\}.
\]
For $V+\nu>0$, the best coefficient is $\lambda^*=V/(V+\nu)$, with improvement $(1-\phi^h)^2V^2/(V+\nu)$. If $V>0$, improvement occurs exactly for $0<\lambda<2V/(V+\nu)$.
\end{proposition}
\begin{proof}
Write the future change as $-(1-\phi^h)z_t+u_{t,h}$, with $u_{t,h}$ centered and independent of $z_t,\xi_t$. Subtract the random-walk squared error from that of the candidate. Taking expectations cancels $Eu_{t,h}^2$ and all innovation cross terms, leaving
$(1-\phi^h)^2E[(z_t-\lambda v_t)^2-z_t^2]$.
The covariance $E[z_tv_t]=V$ and variance $E[v_t^2]=V+\nu$ give the formula. Minimizing the quadratic and solving its strict-negativity inequality prove the remaining claims.
\end{proof}
Consequently a theoretically relevant predictor can worsen forecasting when used without sufficient shrinkage. Parameter learning and predictor revisions are not free: any operational rule must compute its coefficient and predictor using information available at the forecast date.

\section*{A delayed, overlap-aware comparison}
Fix a finite pool of $M$ prespecified currency--horizon--algorithm comparisons, each with real-time forecasts. For one horizon, let $g_t=\widehat s_{t+h}-s_t$ and $U_t=s_{t+h}-s_t$. Assume the declared evaluation class satisfies $|g_t|,|U_t|\le B$ almost surely; these are bounded-law restrictions, not claims about unrestricted exchange rates. The realized loss difference is
$d_t=g_t^2-2g_tU_t$, so $|d_t|\le3B^2$. Put $\mu_t=E[d_t\mid\mathcal I_t]$ and $\overline\mu=T^{-1}\sum_t\mu_t$. Labels arrive after $h$ periods and are available before the next forecast at that same date.

\begin{theorem}[An honest interval for average conditional loss]
Let $k=\min(h,T)$. Simultaneously for the $M$ predeclared comparisons, with probability at least $1-\alpha$,
\[
 |\overline d-\overline\mu|
 \le6B^2\sqrt{\frac{k\log(2kM/\alpha)}{2T}}.
\]
The event holds for adaptive real-time forecasting rules and overlapping forecast errors under the stated conditional boundedness, without independence across consecutive dates.
\end{theorem}
\begin{proof}
Partition dates into $k$ residue subsequences. Along each, $d_t-\mu_t$ is a martingale difference: $g_t$ is measurable at prediction and the preceding label in that subsequence is then observed. Its conditional range has length at most $6B^2$. The conditional Hoeffding exponential bound, iterated along a subsequence of length $n_r$, gives
$|\sum_{t\in r}(d_t-\mu_t)|\le6B^2\sqrt{n_r\log(2kM/\alpha)/2}$
except with probability $\alpha/(kM)$. Union bound across subsequences and comparisons, then use $\sum_r\sqrt{n_r}\le\sqrt{kT}$ and divide by $T$. If comparisons have different $k$, assigning failure $\alpha/M$ separately to each still gives the stated simultaneous conclusion.
\end{proof}
A negative upper endpoint certifies negative average conditional loss on this evaluation path. This is explicitly not a confidence interval for a fixed unconditional stationary mean without an additional stability or sampling argument. Under iid nonoverlapping evaluation blocks, an ordinary bounded-mean interval instead estimates that fixed expectation. Unbounded squared errors require moment or tail assumptions and a different concentration argument. Searching over unregistered algorithms invalidates the finite-pool guarantee unless the enlarged search is also covered.

\section*{Accuracy and trading are different targets}
For a separately specified funded currency or forward payoff $R_{t+1}$ per unit position, let a predictable position $a_t\in[-1,1]$ incur proportional transaction cost $c|a_t|$. Conditional expected net payoff is
$a_tE[R_{t+1}\mid\mathcal I_t]-c|a_t|$. Maximizing this piecewise linear expression gives
$\max\{|E[R_{t+1}\mid\mathcal I_t]|-c,0\}$.
Thus a predictable signal can have a positive squared-forecast gain while all feasible positions in this trading class have nonpositive expected net gain if conditional expected payoffs never exceed costs. Log-rate forecasts, currency excess returns, funding rates and utility risk adjustments require their own mapping; none is silently equated here.

\section*{Remaining gap}
The note proves improvement under a supplied predictable-component model and gives an honest restricted evaluation procedure. It does not establish that any particular currency has those stable primitives, learn them optimally across breaks, or demonstrate empirical improvement on untouched observations. Uniform claims across pairs, aggregate gains and implementable trading benefits therefore remain separate unresolved targets.

\begin{thebibliography}{9}
\bibitem{exchange} B. B. Bakker (2024). \emph{Reconciling Random Walks and Predictability: A Dual-Component Model of Exchange Rate Dynamics}. IMF Working Paper 24/252, December 2024, abstract and Sections 9--10. \url{https://www.imf.org/-/media/files/publications/wp/2024/english/wpiea2024252-print-pdf.pdf}.
\end{thebibliography}
\end{document}
